\documentclass[11pt]{ctexart}
\usepackage[a4paper,margin=1in]{geometry}
\usepackage{graphicx}
\usepackage{xcolor}
\usepackage{hyperref}
\definecolor{refgray}{gray}{0.45}
\title{\textbf{市场最新观点汇总}\\[0.4em]{\large AI叙事盖帽宏观焦虑，但通胀粘性与利率重置正在重塑资产定价逻辑}}
\author{KC桌面 自动生成}
\date{260519}
\begin{document}
\maketitle
\section*{一页摘要}
\begin{itemize}
  \item 亚太市场反弹由AI叙事驱动，而非宏观基本面全面改善，投资者接受“低增长+高通胀”组合并重新配置仓位。
  \item 全球基金经理现金水平骤降至3.9\%，触发历史卖出信号，市场拥挤度接近极端。
  \item 通胀而非增长正成为驱动利率上行的核心因子，股债负相关性已回到1990年代末水平。
  \item 资产重估驱动力从“增长能见度”转向“资产可替代性”，重资产、低折旧行业获得结构性溢价。
  \item 中国资产叙事从宏观博弈转向产业兑现，AI和能源转型资本开支周期成为出口份额扩张的结构性支撑。
\end{itemize}
\section{宏观与利率：通胀粘性主导，央行“冷火鸡”时代开启}
\textbf{市场定价正从增长驱动转向通胀驱动，全球央行渐进主义退场，期限溢价系统性上升。}

\begin{itemize}
  \item 全球基金经理调查显示，净81\%的受访者预期通胀走高，市场正在定价“滞胀但不衰退”的基准情景，投资者不再等待宏观环境完美，而是选择“带病前行”。
  \item 美股与10年期美债收益率的2个月滚动相关性已降至接近-0.8，与1998-2000年相当，当利率上升由通胀驱动而非增长驱动时，股市消化能力急剧下降。
  \item 全球债券市场出现同步重定价，10年期美债收益率突破4.60\%，长端利率涨幅超过短端，市场不再相信央行能通过渐进操作平滑长期利率，开始为“冷火鸡”式政策突变定价。
  \item 欧洲面临“双经济体”结构性撕裂，能源冲击推高通胀预期至3.1\%，市场共识可能低估了通胀并高估了增长，服务业成为本轮冲击的首要传导渠道。
  \item 中国央行5月通过买断式逆回购净回笼1万亿元中长期流动性，年内累计净投放骤降至2900亿元，显著低于历史均值，流动性从“极松”转向“中性偏紧”。
\end{itemize}
\begin{figure}[htbp]
\centering
\includegraphics[width=0.88\linewidth]{figures/fig_002.jpg}
\caption{Exhibit 2 - Exhibit 2: Growth expectations stay in negative territory, but have eased significantly compared to last month Net \% expecting a stronger Global / APAC ex-Japan economy}
\end{figure}
\begin{figure}[htbp]
\centering
\includegraphics[width=0.88\linewidth]{figures/fig_003.jpg}
\caption{Exhibit 3 - Exhibit 3: Inflation expectations remain elevated, with a net 81\% of respondents anticipating higher inflation over the next 12 months, the highest level since Apr 2022 Net \% expecting higher inflation in Asia Pacific ex-Japan in t}
\end{figure}
\begin{figure}[htbp]
\centering
\includegraphics[width=0.88\linewidth]{figures/fig_029.jpg}
\caption{Exhibit 2 - Exhibit 2: Equity/bond yield correlation has reached the most negative levels since the 1990s 2-month US equity/US 10y bond yield correlation}
\end{figure}
\begin{figure}[htbp]
\centering
\includegraphics[width=0.88\linewidth]{figures/fig_030.jpg}
\caption{Exhibit 3 - Exhibit 3: Inflation has been the driver of more negative equity/bond yield correlation}
\end{figure}
\begin{figure}[htbp]
\centering
\includegraphics[width=0.88\linewidth]{figures/fig_032.jpg}
\caption{Exhibit 5 - Exhibit 5: US 10y yields sold off sharply last week 1-week change in US 10y rate in standard deviation terms}
\end{figure}
\begin{figure}[htbp]
\centering
\includegraphics[width=0.88\linewidth]{figures/fig_268.jpg}
\caption{Figure 4 - Figure 4: Globally traded liquid yield curves - Japan re-joins the peloton}
\end{figure}
\begin{figure}[htbp]
\centering
\includegraphics[width=0.88\linewidth]{figures/fig_269.jpg}
\caption{Figure 5 - Figure 5: 10y US vs Germany Spread moves across the Hormuz timeline}
\end{figure}
{\small\color{refgray}\textbf{References:} [R001] 亚太市场正在经历一场“AI叙事”对“宏观焦虑”的系统性盖帽; [R002] 投资者正在押注一个从未被验证过的剧本; [R003] 欧洲的真正风险不是滞胀，而是“双经济体”结构性撕裂的加速; [R005] 市场真正在定价的不是增长，而是“通胀粘性”下的股债再平衡; [R012] 市场真正低估的不是需求，而是流动性的边际收紧; [R016] 全球债券市场正在发出同一个信号：央行“冷火鸡”时代已经开启}

\section{AI/算力：从GPU到GW，供给端再定价成为核心变量}
\textbf{AI投资超级周期正从训练向推理扩散，计算容量（GW）而非芯片数量成为真正瓶颈，供给侧的硬约束正在重塑设备投资格局。}

\begin{itemize}
  \item 四大超大规模云厂商2024-2027年累计资本开支约2万亿美元，2027年新增约20GW计算容量，相当于AWS前18年总和的四倍，指数级增长节奏表明产能释放进入加速期。
  \item ABF基板2030年供需缺口从15\%扩大至22\%，AI需求从训练向推理、从GPU向ASIC、从单一客户向多元化生态扩散，正在固化供应商的定价权与利润率。
  \item WFE市场增长引擎从单一AI逻辑转向存储与逻辑双轮驱动，2026-2028年上行周期延续，洁净室建设进度与设备产能之间的结构性错配成为真正瓶颈。
  \item 成熟制程景气上行周期启动，AI电源管理芯片需求爆发，预计到2027年下半年成熟制程将出现实质性产能短缺，联电等二线代工厂获得溢出机会。
  \item 英伟达在供给端建立的先发优势正转化为可持续议价权，提前锁定供电外壳、先进制程晶圆产能和DRAM供应，这种资源锁定能力在ASIC竞争对手面前几乎不可复制。
\end{itemize}
\begin{figure}[htbp]
\centering
\includegraphics[width=0.88\linewidth]{figures/fig_001.jpg}
\caption{Exhibit 1 - Exhibit 1: Taiwan is seen as the clearest beneficiary of the next phase of the AI cycle FMS views on beneficiaries of the next phase of the AI cycle}
\end{figure}
\begin{figure}[htbp]
\centering
\includegraphics[width=0.88\linewidth]{figures/fig_296.jpg}
\caption{Exhibit 2 - Exhibit 2: Expecting ABF substrate under-supply gap to widen by 2030 vs our prior estimates}
\end{figure}
\begin{figure}[htbp]
\centering
\includegraphics[width=0.88\linewidth]{figures/fig_297.jpg}
\caption{Exhibit 4 - Exhibit 4: We are entering the AI diffusion era}
\end{figure}
\begin{figure}[htbp]
\centering
\includegraphics[width=0.88\linewidth]{figures/fig_314.jpg}
\caption{Exhibit 1 - Exhibit 1: Our 2026 revisions are led by DRAM/Logic and 2027 led by DRAM/NAND}
\end{figure}
\begin{figure}[htbp]
\centering
\includegraphics[width=0.88\linewidth]{figures/fig_315.jpg}
\caption{Exhibit 3 - Exhibit 3: Our 2027 DRAM WFE forecast of \textbackslash{}\$57bn supplies 31\% bit growth, below where we think unconstrained bit demand is (40\%+)}
\end{figure}
\begin{figure}[htbp]
\centering
\includegraphics[width=0.88\linewidth]{figures/fig_345.jpg}
\caption{Exhibit 1 - Exhibit 1: Greater China foundry capacity could grow in high single digits Y/Y in 2026}
\end{figure}
\begin{figure}[htbp]
\centering
\includegraphics[width=0.88\linewidth]{figures/fig_351.jpg}
\caption{Exhibit 7 - Exhibit 7: MS estimates power semi content in Feynman to reach US\textbackslash{}\$191/kW. Total power semi content (\$)}
\end{figure}
{\small\color{refgray}\textbf{References:} [R011] 2万亿美元资本开支的真正赌注，不是GPU，是GW; [R013] 市场真正低估的不是需求，而是供给侧的再定价; [R020] ABF基板供需缺口正在被AI重新定价，而非被产能扩张解决; [R021] 市场真正低估的不是需求，而是供给侧的硬约束正在重塑设备投资格局; [R022] 市场真正低估的不是AI GPU，而是成熟制程的供给缺口; [R038] 市场真正低估的不是需求，而是英伟达在供给侧的定价权}

\section{能源与大宗：地缘冲击加速结构性分化，数据中心需求重塑商品定价}
\textbf{伊朗冲突加速全球经济内部分化，数据中心建设为资源品注入新的结构性增长动力，市场对供给端再定价能力认知不足。}

\begin{itemize}
  \item 全球MAP指标自伊朗冲突以来持续录得下行意外，5月触及-3.5，实际经济数据系统性低于预期，市场对冲突后经济韧性的判断可能过于乐观。
  \item 欧元区4月CAI仅+0.8\%，德国跌至-0.6\%，而印度跳升至+7.9\%，中国骤降2.3个百分点，地缘冲击正在加速全球经济内部的分化。
  \item 数据中心建设对铜和铁矿石的需求增量正在改变必和必拓等资源企业的商品组合定价逻辑，资产进一步变现的可能性可能带来估值重估。
  \item 欧洲能源成本永久性上移，TTF天然气价格预计2026年四季度仍比冲突前高120\%，化肥价格上涨约30\%正在向食品价格传导，三重结构性成本叠加。
  \item 中国出口份额持续扩张，到2030年可能达到16.5\%，AI超级周期和能源转型两轮全球资本开支周期叠加，中国成为核心供应链节点而非仅需求端受益者。
\end{itemize}
\begin{figure}[htbp]
\centering
\includegraphics[width=0.88\linewidth]{figures/fig_174.jpg}
\caption{Exhibit 1 - Exhibit 1: Our MAP Indicator Has Surprised to the Downside Since the Start of The War in Iran}
\end{figure}
\begin{figure}[htbp]
\centering
\includegraphics[width=0.88\linewidth]{figures/fig_175.jpg}
\caption{Exhibit 2 - Exhibit 2: The Global ex Russia FCI Tightened +5.5bps Last Week Primarily on Long Rates Exhibit 3: Higher 2026 Growth in the UK}
\end{figure}
\begin{figure}[htbp]
\centering
\includegraphics[width=0.88\linewidth]{figures/fig_177.jpg}
\caption{Exhibit 4 - Exhibit 5: Our Global CAI Currently Stands Above Potential}
\end{figure}
\begin{figure}[htbp]
\centering
\includegraphics[width=0.88\linewidth]{figures/fig_215.jpg}
\caption{Figure 1 - Figure 1: ECB wage trackers shows c2.3\% wage growth but UBS forecast CPI inflation to rise to 3.4\%}
\end{figure}
\begin{figure}[htbp]
\centering
\includegraphics[width=0.88\linewidth]{figures/fig_216.jpg}
\caption{Figure 2 - Figure 2: European real wage growth on the Indeed Survey is now negative}
\end{figure}
\begin{figure}[htbp]
\centering
\includegraphics[width=0.88\linewidth]{figures/fig_435.jpg}
\caption{Figure 4 - Figure 4: We expect overseas revenue exposure of A-shares to reach 25\% in 2030E}
\end{figure}
{\small\color{refgray}\textbf{References:} [R003] 欧洲的真正风险不是滞胀，而是“双经济体”结构性撕裂的加速; [R008] 伊朗局势正在改写全球经济的短期叙事，但市场低估了结构性分化的深度; [R015] 全球消费股的机会不在整体，而在被错误定价的角落; [R028] 中国出海企业的真实价值不在营收占比，而在利润结构的重构; [R042] 市场真正低估的不是AI需求，而是供给端的再定价能力}

\section{地产与消费：供给端主动收缩，结构性分化加速}
\textbf{地产成交量反弹掩盖领先指标走弱，消费疲软从可选蔓延至必选，但龙头企业通过供给端调整实现盈利韧性。}

\begin{itemize}
  \item 中国地产第20周新房成交环比增长21\%，但领先指标全面走弱：新房搜索量环比下降4\%，二手房带看和认购量分别下降5\%和4\%，价格预期重新转弱。
  \item 四月社零同比仅增0.2\%，补贴退潮后电商增速骤降至0.2\%，家电和手机品类受高基数和补贴效果减弱拖累最为明显。
  \item 白酒行业供给端出现过去十年罕见的主动收缩，部分企业暂停产能建设、控制基酒产量，飞天茅台价格韧性更多来自渠道库存极低和供给收紧，而非需求强劲。
  \item 乳制品是当前最清晰的供给端改善样本，液态奶已率先走出价格战泥潭，蒙牛和伊利均预期从中小厂商手中持续获取份额。
  \item 新能源车市场分化加剧，蔚来和理想周订单环比飙升113\%，而比亚迪、小鹏等品牌环比下滑，新车型集中上市成为检验品牌产品力转化效率的试金石。
\end{itemize}
\begin{figure}[htbp]
\centering
\includegraphics[width=0.88\linewidth]{figures/fig_282.jpg}
\caption{Exhibit 1 - Exhibit 1: Primary GFA sold last week was \$+21\textbackslash{}\%\$ wow and \$+11\textbackslash{}\%\$ yoy in c.75 cities}
\end{figure}
\begin{figure}[htbp]
\centering
\includegraphics[width=0.88\linewidth]{figures/fig_283.jpg}
\caption{Exhibit 2 - Exhibit 2: Primary GFA sold YTD on average was -15\% yoy in c.75 cities, and -15\%/-50\% vs. 2024/2023 level}
\end{figure}
\begin{figure}[htbp]
\centering
\includegraphics[width=0.88\linewidth]{figures/fig_286.jpg}
\caption{Exhibit 5 - Exhibit 5: Average CSI was -1.0pp wow and +3.9pp yoy Weekly Centraline Salesman Index (CSI) tracker in 5 cities}
\end{figure}
\begin{figure}[htbp]
\centering
\includegraphics[width=0.88\linewidth]{figures/fig_287.jpg}
\caption{Exhibit 6 - Exhibit 6: Average CAI was -0.2pp wow and -3.5pp yoy Weekly Centraline Seller Asking Index (CAI) tracker in 6 cities}
\end{figure}
\begin{figure}[htbp]
\centering
\includegraphics[width=0.88\linewidth]{figures/fig_563.jpg}
\caption{Exhibit 1 - Exhibit 1: I-Moutai active users surged from Jan 1st 2026 when Feitian was officially launched on i-Moutai}
\end{figure}
\begin{figure}[htbp]
\centering
\includegraphics[width=0.88\linewidth]{figures/fig_564.jpg}
\caption{Exhibit 2 - Exhibit 2: \$53\textbackslash{}\%\$ Feitian Moutai product prices Most recent data points as of May 17, 2026. Exhibit 3: 52\% Common Wuliangye product prices}
\end{figure}
{\small\color{refgray}\textbf{References:} [R014] 四月数据的真正信号：政策收紧的“用力过猛”与即将到来的转向; [R018] 市场真正低估的不是成交量反弹，而是供给侧的再定价; [R019] 新车潮不是脉冲，是格局分化的加速器; [R025] 4月电商增速骤降至0.2\%，市场低估的不是需求，而是补贴退潮后的结构性重塑; [R026] 消费数据走弱，但市场真正低估的不是需求，而是供给侧的再定价; [R036] 白酒的底在供给侧，但需求侧的春天还没来; [R037] 市场真正低估的不是需求，而是供给侧的主动收缩; [R039] 消费复苏的真正信号不在总量，而在龙头企业的“供给端再定价”}

\section{区域市场：北亚AI叙事巩固，印度面临“AI缺失”折价}
\textbf{亚太市场反弹由AI叙事驱动，日本、台湾、韩国巩固领先地位，印度因缺乏AI标的成为最可能被减持的市场。}

\begin{itemize}
  \item MSCI亚太指数从3月低点反弹17.6\%创历史新高，驱动力来自AI叙事而非宏观基本面改善，台湾被43\%的受访者视为AI下一阶段最清晰受益者。
  \item 日本12个月预期回报创历史新高至6.9\%，67\%的投资者预期日本央行6月加息，盈利已成为日本股市前景的主要驱动力。
  \item 印度被投资者视为“如果全球增长进一步走弱，最先被减持的市场”，缺乏AI标的是首要担忧，与北亚市场的AI叙事形成鲜明对比。
  \item 亚太（除日本）出现板块轮动，从能源和必需消费转向可选消费、房地产和软件，半导体和银行继续成为日本市场最拥挤的板块。
  \item 亚洲银行AI准备度最高的五大市场享有约35\%的相对ROE溢价，但这更多反映“安全交易”逻辑而非市场对AI效率的远见性定价。
\end{itemize}
\begin{figure}[htbp]
\centering
\includegraphics[width=0.88\linewidth]{figures/fig_001.jpg}
\caption{Exhibit 1 - Exhibit 1: Taiwan is seen as the clearest beneficiary of the next phase of the AI cycle FMS views on beneficiaries of the next phase of the AI cycle}
\end{figure}
\begin{figure}[htbp]
\centering
\includegraphics[width=0.88\linewidth]{figures/fig_008.jpg}
\caption{Exhibit 8 - Exhibit 8: Japan growth expectations bounce back strongly in May, though still below the highs seen earlier in the year Net \% expecting a stronger Japanese economy over the next 12 months}
\end{figure}
\begin{figure}[htbp]
\centering
\includegraphics[width=0.88\linewidth]{figures/fig_011.jpg}
\caption{Exhibit 11 - Exhibit 11: Japan return expectations have climbed to an all-time high of 6.9\% FMS views on expected upside for Japan equities over the next 12 months}
\end{figure}
\begin{figure}[htbp]
\centering
\includegraphics[width=0.88\linewidth]{figures/fig_018.jpg}
\caption{Exhibit 18 - Exhibit 18: The absence of a clear AI play is cited as the primary concern for the Indian market What is your key concern for Indian market?}
\end{figure}
\begin{figure}[htbp]
\centering
\includegraphics[width=0.88\linewidth]{figures/fig_019.jpg}
\caption{Exhibit 19 - Exhibit 19: Japan, Taiwan and Korea consolidated their leadership even further, reinforcing a more concentrated North Asia overweight Asia Pacific market sentiment: Net \% overweight Asia Pacific market sentiment: Net \% overweight (\%}
\end{figure}
\begin{figure}[htbp]
\centering
\includegraphics[width=0.88\linewidth]{figures/fig_020.jpg}
\caption{Exhibit 20 - Exhibit 20: India is viewed as the first one to be reduced across APAC if global growth deteriorates further, in line with its largest underweight positioning If global growth weakens further, you would first reduce exposure to...?}
\end{figure}
\begin{figure}[htbp]
\centering
\includegraphics[width=0.88\linewidth]{figures/fig_022.jpg}
\caption{Exhibit 22 - Exhibit 22: May saw rotation out of Energy and Consumer Staples into Consumer Discretionary (ex Retailing), Real Estate, and Software APAC ex-Japan sector sentiment: MoM ppt change in FMS investor positioning}
\end{figure}
{\small\color{refgray}\textbf{References:} [R001] 亚太市场正在经历一场“AI叙事”对“宏观焦虑”的系统性盖帽; [R027] 亚洲银行AI溢价：一次美丽的巧合，而非市场的远见; [R029] 亚洲银行真正的拐点不是AI，而是成本效率红利已近尾声}

\section{风险偏好与拥挤交易：现金水平触发卖出信号，动量因子过度拥挤}
\textbf{市场情绪从悲观切换到乐观，但现金缓冲消失，上涨集中度接近极端，利率对资产定价的深层重置尚未被充分理解。}

\begin{itemize}
  \item 全球基金经理现金水平从4.3\%骤降至3.9\%，触发历史卖出信号，历史上24次类似信号后全球股票中位数四周跌幅为1\%。
  \item 标普500今年迄今约10\%的回报中，TMT板块贡献了85\%，盈利修正驱动力高度集中在科技和能源两大板块，上涨根基并非广泛的盈利复苏。
  \item “做多全球半导体”连续被评为最拥挤交易（73\%），动量因子相对市场比值飙升至115，远高于2025年中基准线，一旦强势股基本面松动，回撤压力将高度集中。
  \item 风险偏好指标飙升至1991年以来第99百分位，美国散户交易量自4月中旬暴增28\%，当所有人都已上车，后续买方力量主要来自稳定性和持续性较弱的资金。
  \item 投资者同时押注经济“不着陆”、利率下行和盈利大幅改善，这个组合在历史上从未被验证过，若后续经济数据不及预期而仓位已打满，调整速度可能比历史上任何一次都快。
\end{itemize}
\begin{figure}[htbp]
\centering
\includegraphics[width=0.88\linewidth]{figures/fig_158.jpg}
\caption{Exhibit 1 - Exhibit 1: Year-to-date EPS revisions have been strong for IT and Energy - the two main drivers of market returns MSCI AC World sectors and Global Regions. Local currency}
\end{figure}
\begin{figure}[htbp]
\centering
\includegraphics[width=0.88\linewidth]{figures/fig_160.jpg}
\caption{Exhibit 3 - Exhibit 3: Momentum vs. Market MSCI Momentum Indices relative price return (USD)}
\end{figure}
\begin{figure}[htbp]
\centering
\includegraphics[width=0.88\linewidth]{figures/fig_161.jpg}
\caption{Exhibit 4 - Exhibit 4: The rebound in technology has reflected the valuation opportunity following the de-rating earlier in the year Distribution of returns of World Tech vs. World ex. TMT, data since 1973}
\end{figure}
\begin{figure}[htbp]
\centering
\includegraphics[width=0.88\linewidth]{figures/fig_162.jpg}
\caption{Exhibit 5 - Exhibit 5: Equity market valuations have increased and largely shrugged off the rise in bond yields, pushing down equity risk premia Global market implied ERP (\%)}
\end{figure}
\begin{figure}[htbp]
\centering
\includegraphics[width=0.88\linewidth]{figures/fig_163.jpg}
\caption{Exhibit 6 - Exhibit 6: The correlation of equities to bond yields has turned negative 3m rolling correlation of daily changes}
\end{figure}
\begin{figure}[htbp]
\centering
\includegraphics[width=0.88\linewidth]{figures/fig_164.jpg}
\caption{Exhibit 7 - Exhibit 7: Sharp bond yield moves have coincided with negative equity returns Average global equities returns depending on absolute moves in US 10-year yields (data since 2000)}
\end{figure}
{\small\color{refgray}\textbf{References:} [R002] 投资者正在押注一个从未被验证过的剧本; [R005] 市场真正在定价的不是增长，而是“通胀粘性”下的股债再平衡; [R007] 市场真正低估的不是动能，而是利率对资产定价的深层重置}

\section*{结语}
综合来看，市场正处于AI叙事主导的结构性分化期，但通胀粘性、利率重置和仓位拥挤三重风险正在累积。真正的机会不在于追逐共识，而在于识别供给端再定价、资产重估驱动力切换以及区域结构性分化中被低估的变量。
\vfill
{\small\color{refgray}Personal reading notes and learning share only. Not investment advice.}
\end{document}
